% ===== PRÉAMBULE CANONIQUE — à coller VERBATIM en tête de CHAQUE .tex =====
\documentclass[11pt,a4paper]{article}
\usepackage[utf8]{inputenc}\usepackage[T1]{fontenc}\usepackage{lmodern}
\usepackage[french]{babel}
\usepackage{amsmath,amssymb,mathtools}\usepackage{icomma}
\usepackage{siunitx}\sisetup{locale=FR}  % \qty{}{}, \si{}, \ang{}
\usepackage{xcolor}\usepackage{colortbl}
\usepackage{tikz}\usepackage{pgfplots}\pgfplotsset{compat=1.18}
\usepackage[siunitx,europeanresistors]{circuitikz} % circuits (élec.)
\usepackage[most]{tcolorbox}\usepackage{enumitem}\usepackage{array,booktabs}
\usepackage[a4paper,margin=2.2cm]{geometry}
\usetikzlibrary{arrows.meta,calc,angles,quotes,positioning,patterns,%
  backgrounds,shapes.geometric,decorations.pathmorphing,decorations.markings,intersections}
\definecolor{primary}{RGB}{222,49,99}\definecolor{primarydark}{RGB}{150,22,55}
\definecolor{defcol}{RGB}{0,114,178}\definecolor{methcol}{RGB}{52,92,148}
\definecolor{excol}{RGB}{0,135,90}\definecolor{attncol}{RGB}{200,40,55}
\tcbset{boxbase/.style={enhanced,breakable,boxrule=0.9pt,arc=2.5pt,
  left=3mm,right=3mm,top=2.3mm,bottom=2.3mm,fonttitle=\bfseries,coltitle=white}}
% --- boîtes TUTORIEL (noms EXACTS, ne pas renommer) ---
\newtcolorbox{cle}[1][]{boxbase,colback=defcol!6,colframe=defcol,
  title={\textsc{À retenir}\ifx\relax#1\relax\relax\else\textmd{ : \itshape #1}\fi}}
\newtcolorbox{methode}[1][]{boxbase,colback=methcol!7,colframe=methcol,sharp corners,
  title={\textsc{Méthode}\ifx\relax#1\relax\relax\else\textmd{ --- \itshape #1}\fi}}
\newtcolorbox{exemplebox}[1][]{boxbase,colback=excol!5,colframe=excol,
  title={\textsc{Exemple}\ifx\relax#1\relax\relax\else\textmd{ : \itshape #1}\fi}}
\newtcolorbox{piege}[1][]{boxbase,colback=attncol!6,colframe=attncol,
  title={\textsc{Piège}\ifx\relax#1\relax\relax\else\textmd{ : \itshape #1}\fi}}
\newtcolorbox{remarque}[1][]{boxbase,colback=black!4,colframe=black!45,
  title={\textsc{Remarque}\ifx\relax#1\relax\relax\else\textmd{ : \itshape #1}\fi}}
% --- boîtes EXERCICE / CORRIGÉ (noms EXACTS, auto-comptées) ---
\newtcolorbox[auto counter]{exo}[1][]{boxbase,colback=primary!4,colframe=primary,
  title={\textsc{Exercice}~\thetcbcounter\ifx\relax#1\relax\relax\else\textmd{ --- \itshape #1}\fi}}
\newtcolorbox[auto counter]{corrige}[1][]{boxbase,colback=excol!4,colframe=excol,
  title={\textsc{Corrigé}~\thetcbcounter\ifx\relax#1\relax\relax\else\textmd{ --- \itshape #1}\fi}}
\newcommand{\vv}[1]{\vec{#1}}\newcommand{\ds}{\displaystyle}
\newcommand{\R}{\mathbb{R}}\newcommand{\N}{\mathbb{N}}\newcommand{\Z}{\mathbb{Z}}
% (un \usepackage{fancyhdr}+en-tête est possible ; il est ignoré côté web)
% =========================================================================
\begin{document}
\sisetup{per-mode=symbol}

\begin{corrige}[tir oblique complet]
\begin{enumerate}[label=\alph*)]
  \item $v_{0x}=20\cos\ang{30}=10\sqrt3\approx\qty{17,3}{\metre\per\second}$ ; $v_{0y}=20\sin\ang{30}=\qty{10}{\metre\per\second}$.
  \item $t_s=\dfrac{v_{0y}}{g}=\dfrac{10}{10}=\qty{1}{\second}$ ; $T=2t_s=\qty{2}{\second}$.
  \item $y_{\max}=\dfrac{v_{0y}^2}{2g}=\dfrac{100}{20}=\qty{5}{\metre}$ ; $x_p=v_{0x}\,T=10\sqrt3\times2=20\sqrt3\approx\qty{34,6}{\metre}$.
\end{enumerate}
\end{corrige}

\begin{corrige}[la vitesse au sommet]
\begin{enumerate}[label=\alph*)]
  \item $v_{0x}=25\times0{,}6=\qty{15}{\metre\per\second}$ ; $v_{0y}=25\times0{,}8=\qty{20}{\metre\per\second}$.
  \item Au sommet, $v_y=0$ : il ne reste que la composante horizontale, $v=v_{0x}=\qty{15}{\metre\per\second}$.
  \item $y_{\max}=\dfrac{v_{0y}^2}{2g}=\dfrac{20^2}{20}=\qty{20}{\metre}$.
\end{enumerate}
\end{corrige}

\begin{corrige}[deux angles complémentaires]
\begin{enumerate}[label=\alph*)]
  \item $x_p(\ang{30})=\dfrac{20^2\sin\ang{60}}{10}=40\times\dfrac{\sqrt3}{2}=20\sqrt3\approx\qty{34,6}{\metre}$ ;
        $x_p(\ang{60})=\dfrac{20^2\sin\ang{120}}{10}=20\sqrt3\approx\qty{34,6}{\metre}$. \textbf{Portées identiques} (angles complémentaires).
  \item $y_{\max}(\ang{30})=\dfrac{(20\sin\ang{30})^2}{20}=\dfrac{10^2}{20}=\qty{5}{\metre}$ ;
        $y_{\max}(\ang{60})=\dfrac{(20\sin\ang{60})^2}{20}=\dfrac{(10\sqrt3)^2}{20}=\dfrac{300}{20}=\qty{15}{\metre}$. Le tir à \ang{60} monte \textbf{trois fois plus haut}.
\end{enumerate}
\end{corrige}

\begin{corrige}[retrouver l'angle de tir]
\begin{enumerate}[label=\alph*)]
  \item $x_p=\dfrac{v_0^2\sin(2\theta)}{g}\Rightarrow \sin(2\theta)=\dfrac{x_p\,g}{v_0^2}=\dfrac{20\times10}{400}=0{,}5$.
  \item $2\theta=\ang{30}$ ou $2\theta=\ang{150}$, soit $\theta=\ang{15}$ ou $\theta=\ang{75}$ (deux angles complémentaires).
\end{enumerate}
\end{corrige}

\begin{corrige}[tir à 45° depuis le sol]
\begin{enumerate}[label=\alph*)]
  \item $y_{\max}=\dfrac{v_0^2\sin^2\ang{45}}{2g}=\dfrac{1600\times0{,}5}{20}=\qty{40}{\metre}$.
  \item $x_p=\dfrac{v_0^2\sin\ang{90}}{g}=\dfrac{1600\times1}{10}=\qty{160}{\metre}$.
\end{enumerate}
\end{corrige}

\end{document}
